\chapter*{List of Publications}
\addcontentsline{toc}{chapter}{List of Publications}

The following publications have been derived from the research work presented in this thesis:

\cite{ref1} Author Name, Supervisor Name, and Co-Author Name, ``Hybrid Adaptive Transformer Differential Protection Using a DWT-LSTM Intelligent Framework with Enhanced CT Saturation and Noise Robustness,'' IEEE Access, vol. XX, pp. XXXXX--XXXXX, 2025. DOI: 10.1109/ACCESS.2025.XXXXXXX.

Summary: This journal paper presents the complete DWT-LSTM framework for power transformer differential protection, including the three-level db4 wavelet-decomposition methodology, the compact six-dimensional energy feature-extraction pipeline, and the LSTM classification architecture with global temporal attention. The paper reports an overall classification accuracy of 98.89\%, with 100\% dependability and 100\% security in the hybrid configuration and a sub-cycle operating time below 17.2 ms for internal faults. Comprehensive evaluation under CT saturation, measurement noise, high-impedance faults, and varying magnetizing conditions is provided, along with comparative analysis against seven existing methods.

\cite{ref2} Author Name, Supervisor Name, and Co-Author Name, ``Multi-Resolution Feature Extraction for Transformer Protection: A Discrete Wavelet Transform Approach with Mutual Information-Based Feature Selection,'' Proceedings of the IEEE International Conference on Power Systems (ICPS), pp. XXX--XXX, 2025. DOI: 10.1109/ICPS.XXXX.XXXXXXX.

Summary: This conference paper focuses on the DWT-based feature-extraction methodology. It presents a systematic analysis of mother-wavelet selection (comparing db4, db6, sym4, sym5, coif3, and bior2.2), decomposition-level optimization, and energy feature extraction from wavelet coefficients. A mutual-information-based feature-importance ranking is introduced, demonstrating that the six energy features---with the three detail-energy components capturing 72.3\% of the discriminative information---are sufficient for reliable four-class discrimination. Ablation-study results on feature-subset size are also presented.

\cite{ref3} Author Name, Supervisor Name, and Co-Author Name, ``Comparative Analysis of Intelligent Methods for Transformer Differential Protection Under Challenging Operating Conditions,'' International Journal of Electrical Power \& Energy Systems, vol. XXX, pp. XXXXX--XXXXX, 2025. DOI: 10.1016/j.ijepes.2025.XXXXXX.

Summary: This journal paper provides a comprehensive comparative analysis of seven transformer differential protection methods, including conventional harmonic restraint, wavelet-ANN, wavelet-SVM, wavelet-PNN, wavelet-CNN, and deep-learning approaches. The comparison is conducted under a unified simulation framework covering five challenging operating conditions: CT saturation, measurement noise, high-impedance faults, severe inrush conditions, and combined worst-case scenarios. The proposed DWT-LSTM framework is shown to outperform all compared methods across the evaluated metrics.

